\documentclass[11pt,a4paper]{article}
\usepackage[margin=2.4cm]{geometry}
\usepackage[T1]{fontenc}
\usepackage[utf8]{inputenc}
\usepackage{amsmath,amssymb,amsthm}
\usepackage{booktabs}
\usepackage{array}
\usepackage{longtable}
\usepackage{enumitem}
\usepackage{tikz}
\usepackage{hyperref}
\hypersetup{colorlinks=true,linkcolor=blue!50!black,urlcolor=blue!50!black}

\newcommand{\Z}{\mathbb{Z}}
\newcommand{\R}{\mathbb{R}}
\newcommand{\Ztw}{\Z_{12}}
\newcommand{\PC}{\Z_{12}}
\newcommand{\lean}[1]{\texttt{#1}}
\newcommand{\boxrow}[1]{\texttt{#1}}

\title{\bf Simple common ground:\\
\emph{The Topos of Music}, Kalkulus I and Lineáris algebra I\\[2mm]
\large and what to actually play on the piano and the drums}
\author{Companion to \lean{RequestProject/MusicOverlap.lean}}
\date{}

\begin{document}
\maketitle

\begin{abstract}
Mazzola's \emph{The Topos of Music} is a big book, but everything in it is built
on a handful of objects that are already \emph{Tárgytematika} elements of your
two courses: a cyclic group, its affine maps, a group action and its invariants,
a logarithm, and a product of two cyclic groups. This note lists those overlaps
one by one, says which exercise of \lean{kalkulus\_gyakorlo.pdf} is the same
mathematics, states the corresponding musical fact, and then --- the point of
the exercise --- gives the simplest pattern that realises it at the keyboard or
on the kit. Every mathematical claim made here is proved in the Lean file
\lean{RequestProject/MusicOverlap.lean}; the theorem name is printed next to the
claim, and Section~\ref{sec:index} is an index of them.
\end{abstract}

\section*{The answer in one table}

\begin{longtable}{@{}p{3.1cm}p{3.9cm}p{3.6cm}p{4.4cm}@{}}
\toprule
\textbf{Shared object} & \textbf{Syllabus item} & \textbf{Mazzola} & \textbf{What you play}\\
\midrule
\endhead
$\log$ / $\exp$, inverse pair
 & Kalkulus: \emph{exponenciális függvények és inverzeik}; gy.\ 2.1e, 2.3, 2.13c
 & pitch is $\log$ of frequency; equal temperament $2^{1/12}$
 & tune a string; hear that ``the same interval'' means ``the same ratio''\\
Quotient $\Z\to\Ztw$, periodicity
 & Kalkulus: periodic functions, gy.\ 2.16g; Lin.\ alg.: homomorphism, kernel
 & octave equivalence, pitch classes
 & play a chord in any octave: same chord\\
Affine maps $x\mapsto \pm x+n$
 & Kalkulus gy.\ 2.16 (\emph{elemi transzformációs lépések}); gy.\ 2.7 ($f^{-1}=f$)
 & transposition $T_n$, inversion $I_n$; the symmetry group
 & transpose a riff; play a melody in mirror image\\
Group, composition, inverse
 & Lin.\ alg.: \emph{mátrixok inverze}; Kalkulus gy.\ 2.18 (composition)
 & the $T/I$ group, $24$ elements (dihedral)
 & the $12$ keys and their mirrors: one riff, $24$ versions\\
Group action, orbit, invariant
 & Lin.\ alg.: linear maps, invariants, determinant
 & interval content of a chord
 & why major and minor are the same chord, upside down\\
Cyclic group again, but for time
 & Lin.\ alg.: $\Z_n$; Kalkulus: periodicity
 & onset sets, rhythm as a subset of $\Z_n$
 & tresillo, clave, Euclidean rhythms on the kit\\
Product / CRT $\Ztw\cong\Z_3\times\Z_4$
 & Lin.\ alg.: direct product; Kalkulus gy.\ 1.4d (index set is a product)
 & the ``time--pitch'' product of coordinates
 & the $3$-against-$4$ polyrhythm\\
\bottomrule
\end{longtable}

\noindent
Everything below is one section per row.

\section{Pitch is a logarithm (Kalkulus: exponential and its inverse)}

Frequencies multiply; keyboards add. The map between the two is
\[
  \operatorname{pitch}(f) \;=\; 69 + 12\log_2\!\left(\frac{f}{440}\right),
\]
so that A4 $=440$\,Hz is note $69$ (\lean{pitch}).

\begin{itemize}[leftmargin=1.4em]
\item \textbf{Octave $=$ doubling.} $\operatorname{pitch}(2f)=\operatorname{pitch}(f)+12$
      (\lean{pitch\_octave}). The multiplicative group $(\R_{>0},\cdot)$ is carried onto
      the additive group $(\R,+)$: this is $\lg(xy)=\lg x+\lg y$, the functional
      equation the Kalkulus course uses all semester.
\item \textbf{Only the ratio matters.}
      $\operatorname{pitch}(f)-\operatorname{pitch}(g)=12\log_2(f/g)$ (\lean{pitch\_sub}).
      A melody transposed to another key is the \emph{same} melody because the ratios
      are unchanged. This is exactly why transposition is a symmetry and not a
      distortion.
\item \textbf{Equal temperament.} $\bigl(2^{1/12}\bigr)^{12}=2$
      (\lean{semitone\_pow\_twelve}), and multiplying a frequency by $2^{1/12}$ raises
      the pitch by one semitone (\lean{pitch\_semitone}).
\end{itemize}

\paragraph{What to play (piano).}
Play C, then the C an octave up: the string vibrates twice as fast. Then play
C--G (ratio $\approx 3/2$, $7$ semitones) and then G--D: the same \emph{sound
relation} twice, i.e.\ the same \emph{added} distance $7+7=14\equiv 2 \pmod{12}$
on the keyboard. You have just performed $\log$.

\paragraph{Why it belongs to the syllabus.}
Gyakorló 2.1e and 2.13c compute with $\lg$; 2.3 and 2.4 compute inverse
functions. Pitch and frequency are an inverse pair of exactly that kind:
$\operatorname{pitch}$ and $f\mapsto 440\cdot 2^{(f-69)/12}$.

\section{Octave equivalence is a quotient homomorphism}

Write $\PC=\Z/12\Z$ for the twelve notes, $0=\mathrm{C}$, $1=\mathrm{C}\sharp$,
\dots, $11=\mathrm{B}$ (\lean{PC}, \lean{C}, \lean{D}, \dots). The single
sentence
\[
  n+12 \equiv n \pmod{12} \qquad (\lean{pc\_octave})
\]
is the whole of ``a C is a C in any octave'', and
$\overline{m+n}=\overline{m}+\overline{n}$ (\lean{pc\_add}) says that the
quotient map preserves intervals.

For the linear algebra course this is a homomorphism with kernel $12\Z$; for the
calculus course it is periodicity, the same thing that makes
$\sin$ and $\cos$ (gyakorló 2.16g, 3.7) behave the way they do.

\paragraph{What to play (piano).}
Take the C major triad C--E--G and re-voice it: E--G--C, then G--C--E. As sets of
pitch classes these are the same three residues; only the octave representatives
changed. Every jazz voicing exercise is a choice of lift along the quotient map
$\Z\to\Ztw$.

\section{Transposition and inversion are the affine maps of 2.16}

Define, on $\PC$,
\[
  T_n(x)=x+n \qquad\text{(transposition)},\qquad
  I_n(x)=n-x \qquad\text{(inversion)} .
\]
These are the maps $x\mapsto \pm x + n$: precisely the ``elemi transzformációs
lépések'' of gyakorló 2.16 (shift the graph, reflect the graph), done over
$\Ztw$ instead of over $\R$.

\begin{itemize}[leftmargin=1.4em]
\item $T_m\circ T_n=T_{m+n}$ \hfill (\lean{T\_comp})
\item $I_n\circ I_n=\mathrm{id}$: inversion is an \emph{involution}
      \hfill (\lean{Inv\_involutive})
\item $I_m\circ I_n=T_{m-n}$: two mirrors make a shift \hfill (\lean{Inv\_comp\_Inv})
\item both are bijections \hfill (\lean{T\_bijective}, \lean{Inv\_bijective})
\end{itemize}

Gyakorló 2.7 asks you to check that $f(x)=\frac{x+1}{x-1}$ satisfies $f^{-1}=f$.
That exercise \emph{is} musical inversion: a self-inverse map. Gyakorló 2.18
(composition of injective maps is injective) is the closure property that makes
these maps a group at all.

\paragraph{The group is the dihedral group of order 24.}
Sending the rotation $r_i$ of the twelve-note clock face to a transposition and
the reflection $sr_i$ to an inversion is an injective group homomorphism
$D_{12}\to \mathrm{Sym}(\PC)$ (\lean{TIhom}, \lean{TIhom\_injective}), so the group
generated by all transpositions and inversions has exactly
\[
  24 = 12\ \text{transpositions} + 12\ \text{inversions}
  \qquad (\lean{TI\_card})
\]
elements. This is the group that music theory calls $T/I$ and that Mazzola treats
as a special case of the affine symmetries of a module --- the same affine
symmetries the linear algebra course meets as $x\mapsto Ax+b$.

\paragraph{The linear algebra behind it: the affine group.}
Write $\operatorname{aff}_{a,b}(x)=ax+b$ (\lean{affine}). Then $T_n=\operatorname{aff}_{1,n}$
and $I_n=\operatorname{aff}_{-1,n}$ (\lean{T\_eq\_affine}, \lean{Inv\_eq\_affine}), affine
maps compose to affine maps (\lean{affine\_comp}), and $\operatorname{aff}_{a,b}$ is
invertible exactly when $a$ is a unit (\lean{affine\_bijective}) --- the modular
version of ``$\det\neq 0$'' from Lineáris algebra I. There are four units mod $12$,
namely $1,5,7,11$ (\lean{units\_card}), so the full affine group has $4\cdot 12=48$
elements and the musically familiar $T/I$ group is the half with $a=\pm1$. The
other two multipliers are the classical ``circle-of-fifths transform'' $x\mapsto 7x$
and its inverse $x\mapsto 5x$: play a chromatic scale but multiply every step by
seven and you are playing the circle of fifths.

\paragraph{What to play (piano), three one-minute drills.}
\begin{enumerate}[leftmargin=1.6em]
\item \textbf{$T$ drill.} Play a two-bar riff on the white keys. Now play it
      starting a minor third higher ($T_3$), then again ($T_6$), then again
      ($T_9$). After four moves you are home: $T_3^4=T_{12}=\mathrm{id}$.
\item \textbf{$I$ drill.} Put your right hand on C. For every note the right hand
      plays $x$ semitones \emph{up}, the left hand plays $x$ semitones
      \emph{down}: the left hand is playing $I_0$. Do it once more with the mirror
      centred on G ($I_7$) and hear the difference.
\item \textbf{Two mirrors $=$ one shift.} Play the riff, invert it around C,
      then invert \emph{that} around G. You are now playing the original riff
      transposed by $7-0=7$ semitones: $I_7\circ I_0=T_7$.
\end{enumerate}

\section{The invariant: interval content}

For a chord $S\subseteq\PC$ let $\operatorname{ivec}(S,i)$ count the ordered pairs of
notes of $S$ lying $i$ semitones apart (\lean{ivec}). Then
\[
  \operatorname{ivec}(T_n S,\,i)=\operatorname{ivec}(S,\,i),\qquad
  \operatorname{ivec}(I_n S,\,i)=\operatorname{ivec}(S,\,-i)
  \qquad (\lean{ivec\_T},\ \lean{ivec\_Inv}).
\]
So interval content is the invariant of the $T/I$ action, in the same sense in
which rank and determinant are invariants of a matrix under change of basis
(Lineáris algebra I). Concretely:
\[
  \{0,4,7\}\ \text{(major)}\ \xrightarrow{\;I_7\;}\ \{0,3,7\}\ \text{(minor)}
  \qquad (\lean{minor\_eq\_inverted\_major}),
\]
so major and minor have the same interval content read backwards
(\lean{triads\_same\_interval\_content}). Two chords built from the same
intervals, two completely different moods: the invariant explains the family
resemblance, the group element explains the difference.

\paragraph{What to play (piano).}
Play C--E--G, then C--E$\flat$--G. Same three gaps $\{3,4,5\}$ in a different
order. Then play the mirror move directly: hold G, and swap E for E$\flat$ ---
one finger, and you have applied $I_7$.

\section{Simple things the group hands you}

All of the following are finite checks, settled in Lean by \lean{decide}.

\begin{itemize}[leftmargin=1.4em]
\item \textbf{The circle of fifths visits all twelve notes}
      (\lean{fifths\_surjective}), because $\gcd(7,12)=1$. Playing
      C--G--D--A--E--B--F$\sharp$--C$\sharp$--G$\sharp$--D$\sharp$--A$\sharp$--F--C
      is a single cycle through the whole keyboard.
\item \textbf{The major scale is seven consecutive fifths}
      (\lean{diatonic\_eq\_seven\_fifths}): stack F--C--G--D--A--E--B, sort, and
      you have C major. The pentatonic scale is five consecutive fifths
      (\lean{pentatonic\_eq\_five\_fifths}): C--D--E--G--A, the shape of the five black keys.
\item \textbf{The major scale has no symmetry}
      (\lean{diatonic\_asymmetric}): shifting it always produces a different set
      of notes, which is why the twelve keys are genuinely twelve keys.
\item \textbf{The whole-tone scale has only two transpositions}
      (\lean{wholeTone\_T2}, \lean{wholeTone\_two\_transpositions}): it is fixed by
      $T_2$. This is Messiaen's ``mode of limited transposition'', and it is the
      reason whole-tone runs sound directionless --- there is no ``home''.
\item \textbf{The augmented triad is an orbit}: C--E--G$\sharp$ is the orbit of C
      under $T_4$ (\lean{augmented\_orbit}), and it too is symmetric.
\end{itemize}

\paragraph{What to play (piano).}
Practise the circle of fifths with the left hand alone, one bar per note.
Then play a whole-tone run up and down and notice that you cannot tell where it
began: symmetric sets carry no information about the tonic, asymmetric sets do.
That single contrast is the musical content of the two theorems above.

\section{Rhythm is the same group, used for time}

A one-bar pattern of $n$ equal slots is a subset of $\Z_n$ (\lean{Pattern}), and
``start the loop somewhere else'' is $x\mapsto x+k$ --- literally the same
operation as transposition (\lean{transposition\_eq\_rotation}). This is
Mazzola's point that pitch and onset are two coordinates of one object, and it is
also the cheapest thing in this whole note to use at the kit.

\paragraph{Euclidean rhythms.} Spread $k$ onsets as evenly as possible over $n$
slots: $E(k,n)=\{\lfloor in/k\rfloor : i<k\}$ (\lean{euclid}). Writing
\texttt{x} for a hit and \texttt{.} for a rest:

\begin{center}
\begin{tabular}{@{}llll@{}}
\toprule
Pattern & Slots & Box notation & Name / where you hear it\\
\midrule
$E(2,4)$ & 4  & \boxrow{x.x.}            & backbeat\\
$E(3,8)$ & 8  & \boxrow{x.x..x..}        & \emph{tresillo} (2--3--3 here) (\lean{euclid\_3\_8})\\
$E(5,8)$ & 8  & \boxrow{xx.xx.x.}        & \emph{cinquillo} (\lean{euclid\_5\_8})\\
$E(5,16)$& 16 & \boxrow{x..x..x..x..x...}& evenly spread clave (\lean{euclid\_5\_16})\\
son clave& 16 & \boxrow{x..x..x...x.x...}& the 3--2 son clave (\lean{clave})\\
$E(4,16)$& 16 & \boxrow{x...x...x...x...}& four-on-the-floor\\
\bottomrule
\end{tabular}
\end{center}

\begin{itemize}[leftmargin=1.4em]
\item \textbf{Asymmetric patterns tell you where beat one is.} The tresillo is
      fixed by no rotation at all (\lean{tresillo\_asymmetric}); rotate it and you
      get a different pattern. That is why it grooves.
\item \textbf{Symmetric patterns do not.} Four-on-the-floor is unchanged by
      rotating one beat (\lean{fourOnFloor\_symmetric}), so on its own it cannot
      tell you where the bar starts --- you need the backbeat or the clave for
      that.
\end{itemize}

\paragraph{What to play (drums), in order of difficulty.}
\begin{enumerate}[leftmargin=1.6em]
\item Kick \boxrow{x...x...x...x...} (four-on-the-floor), snare on \boxrow{....x.......x...}.
      Feel that the kick alone is symmetric and the snare breaks the symmetry.
\item Move the kick to the tresillo \boxrow{x..x..x.} over eight eighths --- count
      \emph{3--3--2} out loud. Keep the snare on beats 2 and 4. This one pattern
      covers a huge amount of pop, reggaetón and Afro-Cuban music.
\item Play the son clave \boxrow{x..x..x...x.x...} on the rim with one hand and
      steady eighths on the hi-hat with the other.
\item Cinquillo \boxrow{xx.xx.x.} with alternating sticking.
\end{enumerate}

\section{Polyrhythm is a product (Chinese remainder)}

Two hands, one hitting every $3$ slots, one every $4$. They coincide exactly
every $12$ slots (\lean{polyrhythm\_3\_4}), and, better,
\[
  \Ztw \;\cong\; \Z_3\times\Z_4 \qquad (\lean{polyIso}),
\]
so a position in the bar \emph{is} the pair (where you are in the three-cycle,
where you are in the four-cycle). This is the universal property of the product
--- the subject of the companion file \lean{RequestProject/ProductOverlap.lean}
--- heard rather than drawn, and $12=3\cdot 4$ independent choices
(\lean{poly\_card}).

\paragraph{What to play.}
Count twelve slots per bar. Right hand on $0,3,6,9$; left hand on $0,4,8$:
\begin{center}
\begin{tabular}{@{}ll@{}}
slots  & \boxrow{0\,1\,2\,3\,4\,5\,6\,7\,8\,9\,10\,11}\\
right (4 hits) & \boxrow{x..x..x..x..}\\
left  (3 hits) & \boxrow{x...x...x...}\\
\end{tabular}
\end{center}
The hands meet only on slot $0$. At the piano, play the same thing as a hemiola:
right hand three notes per bar, left hand four.

\section{What we deliberately left out of Mazzola}

The overlaps above are the elementary layer. \emph{The Topos of Music} goes on to
local and global compositions, denotators, presheaves and Grothendieck topologies,
and its counterpoint model; none of that is needed for the patterns in this note,
and none of it appears in either \emph{Tárgytematika}. What \emph{does} appear in
both, and is used above, is: groups and their actions, homomorphisms and
quotients, affine maps, products, invariants, logarithms and periodicity.

\section{Index of the Lean results}\label{sec:index}

All in \lean{RequestProject/MusicOverlap.lean}, all proved, no \lean{sorry}.

\begin{center}
\begin{tabular}{@{}ll@{}}
\toprule
Lean name & Claim\\
\midrule
\lean{pitch}, \lean{pitch\_octave} & pitch is $\log$ of frequency; octave $=+12$\\
\lean{pitch\_sub}                  & pitch difference depends only on the ratio\\
\lean{semitone\_pow\_twelve}, \lean{pitch\_semitone} & equal temperament\\
\lean{pc\_octave}, \lean{pc\_add}  & octave equivalence is a quotient homomorphism\\
\lean{T}, \lean{Inv}, \lean{T\_comp}, \lean{Inv\_involutive}, \lean{Inv\_comp\_Inv} & the two affine operations\\
\lean{T\_bijective}, \lean{Inv\_bijective} & both are bijections\\
\lean{affine}, \lean{T\_eq\_affine}, \lean{Inv\_eq\_affine}, \lean{affine\_comp} & $T$ and $I$ are the affine maps $x\mapsto ax+b$\\
\lean{affine\_bijective}, \lean{units\_card} & invertible $\iff$ unit multiplier; $48$ affine maps\\
\lean{TIhom}, \lean{TIhom\_injective}, \lean{TI\_card} & the $T/I$ group is dihedral of order $24$\\
\lean{ivec}, \lean{ivec\_T}, \lean{ivec\_Inv} & interval content is the invariant\\
\lean{minor\_eq\_inverted\_major}, \lean{triads\_same\_interval\_content} & major and minor are mirror images\\
\lean{fifths\_surjective} & the circle of fifths is one cycle\\
\lean{diatonic\_eq\_seven\_fifths}, \lean{pentatonic\_eq\_five\_fifths} & scales are stacks of fifths\\
\lean{diatonic\_asymmetric} & the major scale has no transpositional symmetry\\
\lean{wholeTone\_T2}, \lean{wholeTone\_two\_transpositions} & limited transposition\\
\lean{augmented\_orbit} & the augmented triad is a $T_4$-orbit\\
\lean{transposition\_eq\_rotation} & transposing $=$ rotating a loop\\
\lean{euclid}, \lean{euclid\_3\_8}, \lean{euclid\_5\_8}, \lean{euclid\_5\_16} & Euclidean rhythms\\
\lean{tresillo\_asymmetric}, \lean{fourOnFloor\_symmetric} & which grooves point at beat one\\
\lean{polyrhythm\_3\_4}, \lean{polyIso}, \lean{poly\_card} & $3$-against-$4$ is a product\\
\bottomrule
\end{tabular}
\end{center}

\end{document}
